\section{The literal definition}
\label{sec:literal}

This section settles the literal reading of exercise~15(c), in which
every augmenting path has the current shortest level at every round.
Two consecutive rounds of this kind lose at least half of the free
girls, while the Tchoukaillon schedule asks the first round at level~2 to
remove only about a third of them.

\subsection{Uniform consecutive levels force halving}

Throughout this subsection the graph has a perfect matching ($s=M=N$),
and distances are measured in the oriented graph of
Section~\ref{sec:dag}.

\begin{lthm}{U}
Suppose a matching $M'$ with $f$ free girls has $f$ vertex-disjoint SAPs, of
level~$L$ (property~$\Bp$; it holds in particular under property~$\Up$).
Let $M''$ be obtained from $M'$ by augmenting along $a\ge1$
vertex-disjoint SAPs. If every augmenting path for $M''$ has level~$L+1$,
then
\[
   f-a\le a .
\]
In particular, if a round has property~$\Bp$ and the next round has
property~$\Up$ at the next level, the first of the two rounds removes at
least half of the free girls.
\end{lthm}

\begin{proof}
Call the $a$ paths along which we augmented the \emph{selected} paths.
On a selected path $S=g_0\sg b_1\db g_1\sg\cdots\db g_{L-1}\sg b_L$, write
$e_i$ for its edge $g_{i-1}\sg b_i$ outside $M'$; it is a couple of~$M''$.

\emph{Step~1: old distances are still lower bounds.} For $M'$, let $d(v)$
be the distance of $v$ from the free girls, and $h(v)$ the length of a
shortest directed path from $v$ to a free boy. On a selected path these
are exact: $d(g_i)=2i$, $d(b_i)=2i-1$, $h(g_i)=2(L-i)-1$, and
$h(b_i)=2(L-i)$. Passing to $M''$ reverses the arcs of the selected paths
and makes some free vertices matched. The reversed arcs are $b_i\to
g_{i-1}$ and $g_i\to b_i$, and they satisfy $d(\text{head})\le
d(\text{tail})+1$ and $h(\text{tail})\le h(\text{head})+1$: for instance
$d(g_{i-1})=2i-2\le d(b_i)+1$ and $h(b_i)=2(L-i)\le h(g_{i-1})+1$, and
similarly for $g_i\to b_i$. All other arcs are unchanged. So $d$ and $h$
remain lower bounds for the corresponding quantities with respect
to~$M''$.

\emph{Step~2: contacts are aligned.} Let $A$ be an augmenting path for
$M''$ that uses the couple $e_i$, say as its $k$th couple counting from its
free girl. Its part up to $b_i$ has $2k-1$ edges and its part from
$g_{i-1}$ on has $2L+1-2k$ edges, so Step~1 gives $2k-1\ge2i-1$ and
$2L+1-2k\ge2(L-i)+1$. Hence $k=i$: a new augmenting path can use $e_i$ only
as its $i$th couple.

\emph{Step~3: new basic paths meet the selected paths.} Let $P^*$ be the
perfect matching obtained from $M'$ by augmenting along all $f$ of its
disjoint SAPs. The symmetric difference of $M''$ and $P^*$ contains $f-a$
vertex-disjoint augmenting paths for $M''$; call their set $\mathcal A$.
A member of $\mathcal A$ that avoided all selected paths would be a
component of $M'\symdiff P^*$ as well, that is, one of the old SAPs, still
augmenting and of level~$L$, contrary to the hypothesis. So every member
of $\mathcal A$ meets a selected path, and since all vertices of the
selected paths are matched in $M''$, it contains one of their couples
$e_i$.

\emph{Step~4: no selected path meets two members of $\mathcal A$.}
Suppose the selected path $S$ meets two. Among the couples of $S$ that
belong to members of $\mathcal A$, pick two, $e_i\in A_i$ and $e_j\in A_j$
with $i<j$ and $A_i\ne A_j$, such that no couple strictly between them
belongs to a member of $\mathcal A$. Now splice (Fig.~\ref{fig:splice}):
follow $A_j$ from its free girl through $e_j$ to $g_{j-1}$, then walk
backwards along $S$ through $b_{j-1},g_{j-2},\ldots,g_i$ to $b_i$, then
follow $A_i$ from $b_i$ through $e_i$ to its free boy. The three pieces
meet only at their ends, so this is an augmenting path for $M''$. By
Step~2 its number of edges is
\[
  2j\;+\;\bigl(2(j-i)-1\bigr)\;+\;\bigl(2L+2-2i\bigr)\;=\;2L+1+4(j-i)\;>\;2L+1,
\]
so its level exceeds $L+1$, contradicting the hypothesis.

By Steps~3 and~4, each of the $f-a$ members of $\mathcal A$ meets a
selected path, and each of the $a$ selected paths meets at most one of
them. So $f-a\le a$.
\end{proof}

\begin{figure}[ht]
\centering
% Theorem U: two new paths Q_i, Q_j touching the old SAP P at e_i, e_j.
\begin{tikzpicture}[x=1.12cm,y=1cm]
% P after augmentation: g0 b1 g1 b2 g2 b3 g3 b4 g4 b5, couples g_{k-1} b_k
\foreach \lab/\x in {g_0/0,g_1/2,g_2/4,g_3/6,g_4/8}{\node[girl] (\lab) at (\x,0) {$\lab$};}
\foreach \lab/\x in {b_1/1,b_2/3,b_3/5,b_4/7,b_5/9}{\node[boy] (\lab) at (\x,0) {$\lab$};}
% free girls / boys of Q_i and Q_j
\node[girl,free] (qi0) at (1.2,2.0) {};
\node[boy,free]  (qi1) at (1.6,-2.0) {};
\node[girl,free] (qj0) at (5.6,2.0) {};
\node[boy,free]  (qj1) at (5.0,-2.0) {};
\begin{scope}[on background layer]
  % the spliced path
  \draw[sap] (qj0.center) .. controls (6.3,1.4) and (7.2,1.0) .. (b_4.center)
     -- (g_3.center) -- (b_3.center) -- (g_2.center) -- (b_2.center) -- (g_1.center)
     .. controls (1.6,-.8) and (1.2,-1.4) .. (qi1.center);
\end{scope}
% couples e_k of the new matching along P
\draw[couple] (g_0)--(b_1) (g_1)--(b_2) (g_2)--(b_3) (g_3)--(b_4) (g_4)--(b_5);
\draw[edge,line width=.75pt] (b_1)--(g_1) (b_2)--(g_2) (b_3)--(g_3) (b_4)--(g_4);
% Q_i and Q_j (schematic)
\draw[boydraw,line width=.8pt,densely dashed] (qi0) .. controls (2.2,1.5) and (3,1.0) .. (b_2);
\draw[boydraw,line width=.8pt,densely dashed] (g_1) .. controls (1.6,-.8) and (1.2,-1.4) .. (qi1);
\draw[girldraw,line width=.8pt,densely dashed] (qj0) .. controls (6.3,1.4) and (7.2,1.0) .. (b_4);
\draw[girldraw,line width=.8pt,densely dashed] (g_3) .. controls (5.6,-.8) and (5.2,-1.4) .. (qj1);
\node[font=\small,boydraw] at (2.6,1.75) {$A_i$};
\node[font=\small,girldraw] at (7.0,1.75) {$A_j$};
\node[font=\footnotesize,anchor=north] at (2.5,-.32) {$e_i$};
\node[font=\footnotesize,anchor=north] at (6.5,-.32) {$e_j$};
\node[font=\small,anchor=west] at (9.4,0) {$S$};
\end{tikzpicture}

\caption{Step~4 of the proof of Theorem~U. After augmenting along the
selected path $S$, its couples are the edges $e_k$ from $g_{k-1}$ to $b_k$.
If two disjoint augmenting paths $A_i$ and $A_j$ cross $S$ at $e_i$ and
$e_j$, the shaded splice is an augmenting path with $4(j-i)$ more edges
than either of them.}
\label{fig:splice}
\end{figure}

The splice argument also shows that a full set of disjoint new SAPs can
coexist with much longer augmenting paths.

\begin{proposition}
\label{prop:longpaths}
In the situation of Theorem~U, drop the hypothesis that all augmenting
paths for $M''$ have level~$L+1$, and assume only that its SAPs have level
$L+1$. If $M''$ has $q$ vertex-disjoint SAPs, each meeting a selected
path, then $M''$ has an augmenting path with at least
$2L+1+4(\ceil{q/a}-1)$ edges.
\end{proposition}

\begin{proof}
Steps~1 and~2 apply to SAPs. Some selected path $S$ meets at least
$k=\ceil{q/a}$ of the $q$ SAPs. Choose a shortest interval $[i,j]$ of
indices such that the couples $e_i,\ldots,e_j$ of $S$ include contacts with
all $k$ of them. The SAPs owning $e_i$ and $e_j$ have no other contact
inside the interval (otherwise it could be shortened), and $j-i\ge k-1$.
Splice as in Step~4; contacts with other SAPs inside the interval do no
harm, because those paths are not used.
\end{proof}

For example, in the first dag round of Knuth's 16-pair graph three SAPs
are selected, and the next matrix has five disjoint SAPs of level~3, each
meeting a selected one. So the second matching has an augmenting path
with at least $5+4=9$ edges, that is, of level at least~5; the census in
Appendix~\ref{app:knuth16} shows six of them.

\subsection{Consequences}

\begin{corollary}
If property $\Up$ holds in consecutive rounds at levels~$L,L+1,\ldots$,
and the first of them starts with $f>0$ free girls, there are at most
$\floor{\lg f}+1$ such rounds. The bound is attained when $f=1,2,4,8,16,32$,
and we expect it to be attained for every power of~2.
\end{corollary}

\begin{proof}
By Theorem~U every round except the last leaves at most $\floor{f/2}$ free
girls. Appendix~\ref{app:halving} describes ordered graphs, one for each
$f=2^t$, on which Algorithm~H halves the number of free girls in every
round while every augmenting path has the current shortest level
(Fig.~\ref{fig:halving}). For $f\le32$ we verified this by running
Algorithm~H and enumerating all augmenting paths in every round; the
appendix explains why it should hold in general, but we have not written
out a complete proof.
\end{proof}

\begin{figure}[ht]
\centering
\begin{tikzpicture}[]
\draw[pedge] (3.100,1.500) -- (3.100,0.500);
\draw[edge,draw=black!55,line width=.75pt] (3.100,1.500) to[bend left=12] (4.650,0.500);
\draw[pedge] (0.000,1.500) -- (0.000,0.500);
\draw[edge,draw=black!55,line width=.75pt] (0.000,1.500) to[bend left=12] (1.550,0.500);
\draw[pedge] (1.550,1.500) -- (1.550,0.500);
\draw[edge,draw=black!55,line width=.75pt] (1.550,1.500) to[bend left=12] (4.650,0.500);
\draw[pedge] (4.650,1.500) -- (4.650,0.500);
\draw[couple] (0.000,-0.500) -- (0.000,0.500);
\draw[pedge] (0.000,-0.500) -- (0.000,-1.500);
\draw[couple] (1.550,-0.500) -- (1.550,0.500);
\draw[pedge] (1.550,-0.500) -- (1.550,-1.500);
\draw[couple] (3.100,-0.500) -- (3.100,0.500);
\draw[pedge] (3.100,-0.500) -- (3.100,-1.500);
\draw[couple] (4.650,-0.500) -- (4.650,0.500);
\draw[pedge] (4.650,-0.500) -- (4.650,-1.500);
\node[girl,free] at (0.000,1.500) {$u_{1}$};
\node[boy] at (0.000,0.500) {$b_{1}$};
\node[girl] at (0.000,-0.500) {$c_{1}$};
\node[boy,free] at (0.000,-1.500) {$d_{1}$};
\node[girl,free] at (1.550,1.500) {$u_{2}$};
\node[boy] at (1.550,0.500) {$b_{2}$};
\node[girl] at (1.550,-0.500) {$c_{2}$};
\node[boy,free] at (1.550,-1.500) {$d_{2}$};
\node[girl,free] at (3.100,1.500) {$u_{3}$};
\node[boy] at (3.100,0.500) {$b_{3}$};
\node[girl] at (3.100,-0.500) {$c_{3}$};
\node[boy,free] at (3.100,-1.500) {$d_{3}$};
\node[girl,free] at (4.650,1.500) {$u_{4}$};
\node[boy] at (4.650,0.500) {$b_{4}$};
\node[girl] at (4.650,-0.500) {$c_{4}$};
\node[boy,free] at (4.650,-1.500) {$d_{4}$};
\end{tikzpicture}

\caption{The halving graph for $f=4$ ($s=8$), with the matching after~\Hs1.
Free girl~$u_i$ is adjacent to $b_i$ and, if $i<4$, also to $b_{i+2^{\rho(i)}}$,
where $2^{\rho(i)}$ is the largest power of~2 dividing~$i$. Each dag round
merges pairs of neighbouring groups of indices (Appendix~\ref{app:halving}),
and the free girls go $4\to2\to1\to0$. At every stage all augmenting paths
have the same level.}
\label{fig:halving}
\end{figure}

\begin{lthm}[Corollary]{L}
For $s\ge1$, a literal realization of the Tchoukaillon schedule exists
if and only if $s\in\{1,2,3,4,5,8,9\}$.
\end{lthm}

\begin{proof}
The schedule starts with $f=\floor{s/2}$ free girls at level~2 and removes
$a=\ceil{f/3}$ SAPs. If any girls remain, the next round is at level~3 and,
in a literal realization, has property $\Up$; then Theorem~U requires
$f\le2\ceil{f/3}$. This fails for $f=3$ and for every $f\ge5$, which
excludes $s=6,7$ and all $s\ge10$.

For $s=2,4,8$, take the halving graphs with $f=1,2,4$: their free girls go
$1\to0$, $2\to1\to0$, and $4\to2\to1\to0$, which is the schedule in these
cases. For $s=3,5,9$, add one isolated girl--boy edge, numbered last;
\Hs1 matches it and nothing else changes, and $\floor{s/2}$ is unchanged.
For $s=1$ a single edge will do.
\end{proof}

So Knuth's 16-pair graph cannot be repaired into a literal example: no
graph and no input order satisfy the literal definition at $s=16$. Nothing
in this section obstructs the basic-row or operational readings, and
Knuth's graph meets the basic-row reading at $s=16$.
